\chapter{五次三维簇}\label{chap:quintic}
\SourcePage{38}

五次三维簇是人们首次运用镜像对称作出枚举预测的例子 \SourceCite{CdGP}。本章将回顾这一例子，基本沿用 \SourceCite{CdGP} 的方法，同时根据 \SourceCite{Morrison2} 更偏重数学的论述作一些调整。第 \ref{sec:quintic-other-methods} 节将依据 \SourceCite{Givental2} 与 \SourceCite{LLY}，介绍计算同一组枚举数据的另外两种方法。在后面的章节，我们还会多次回到五次三维簇，用它来说明研究镜像对称的不同方法。

\section{五次三维簇的 A 模型关联函数}\label{sec:quintic-a-model}

众所周知，光滑五次超曲面 $V\subset\mathbb P^4$ 是 Calabi--Yau 簇；\hyperref[chap:mirror-constructions]{第 4 章}将建立的理论也会推出这一点。我们考虑与 $(V,\omega)$ 相联系的非线性 $\sigma$ 模型，其中 $\omega=B+iJ$ 是 $V$ 上的复化 Kähler 类。若以 $H$ 表示超平面类，则\hyperref[chap:introduction]{第 1 章}中的 A 模型关联函数~\eqref{eq:a-yukawa} 简化为
\begin{equation}\label{eq:quintic-a}
\langle H,H,H\rangle=5+\sum_{d=1}^{\infty}n_d d^3\frac{q^d}{1-q^d}.
\end{equation}
粗略地说，$n_d$ 是 $V$ 中次数为 $d$ 的有理曲线的数目，而 $q=e^{2\pi i\int_\ell\omega}$，$\ell$ 是 $V$ 中的一条直线。若写 $\omega=tH$，则 $q$ 也可表示为 $e^{2\pi it}$。下一节将看到，$q$ 实际上是紧化 Kähler 模空间上的一个局部坐标。另外，正如\hyperref[chap:introduction]{第 1 章}所解释的，由于这是 A 模型关联函数，式~\eqref{eq:quintic-a} 依赖于 $\omega$，却不依赖于 $V$ 的复结构。

代数几何学家初次见到式~\eqref{eq:quintic-a} 时，觉得它有些奇怪。其中一个原因是，我们甚至还不知道一般五次三维簇对给定次数 $d$ 是否只含有限多条有理曲线。按照 Clemens 猜想，答案应当是肯定的；但迄今只在 $d\leq9$ 的情形下知道该猜想成立 \SourceCite{JK1}。式~\eqref{eq:quintic-a} 却似乎预设它对所有 $d$ 都成立。更令人不安的是，镜像对称竟声称能给出这些曲线的数目 $n_d$ 的显式公式，而且对所有 $d\geq1$ 都适用。\index{Clemens猜想@Clemens 猜想|hyperpage}

\hyperref[chap:gromov-witten]{第~\ref*{chap:gromov-witten} 章}将为式~\eqref{eq:quintic-a} 建立牢固的基础。关键对象是 Gromov--Witten 不变量 $\langle I_{0,3,d}\rangle(H,H,H)$。这个不变量有多种定义方式。从辛几何的观点出发，可以考虑满足适当形变后的 Cauchy--Riemann 方程的 $C^\infty$ 映射 $f:\mathbb P^1\to V$。对 $\mathbb P^1$ 上三个点的像施加合适的条件后，可以证明，在每个同调类中这样的曲线都只有有限多条。于是，$\langle I_{0,3,d}\rangle(H,H,H)$ 被定义为这些曲线按定向所赋符号计数后的总数；这里已不再有来自复结构的优选定向。

\SourcePage{39}
事实上，要把这一想法严格化，有几种不同的做法；\hyperref[sec:gw-symplectic-definition]{第~\ref*{sec:gw-symplectic-definition} 节}将逐一回顾。

\SourceCite{Kontsevich2} 提出了定义 $n_d$ 的另一条途径：使用稳定映射 $f:C\to V$ 的模空间，其中 $C$ 的亏格为零，但不必不可约。由此得到的 Gromov--Witten 不变量定义相当精巧，涉及代数叠和虚拟基本类。\hyperref[sec:gw-algebraic-definition]{第~\ref*{sec:gw-algebraic-definition} 节}将讨论这些主题。

一旦 Gromov--Witten 不变量得到严格定义，我们便可把 A 模型关联函数定义为形式和
\[
\langle H,H,H\rangle=\sum_{d=0}^{\infty}\langle I_{0,3,d}\rangle(H,H,H)q^d.
\]
为弄清它与式~\eqref{eq:quintic-a} 的关系，固定 $d$，设 $k\mid d$。若 $C\subset V$ 是次数为 $k$、法丛为 $\mathcal O_{\mathbb P^1}(-1)\oplus\mathcal O_{\mathbb P^1}(-1)$ 的光滑有理曲线，那么次数为 $d/k$ 的映射 $f:\mathbb P^1\to C$ 就给出被 $\langle I_{0,3,d}\rangle(H,H,H)$ 计入的一类稳定映射。\hyperref[chap:gromov-witten]{第~\ref*{chap:gromov-witten} 章}将看到，这些 $f$ 所组成的族对该不变量贡献一个 $k^3$ 因子。由于每一条次数为 $k\mid d$ 的曲线 $C$ 都如此，因而有理由期待
\begin{equation}\label{eq:multi-cover}
\langle I_{0,3,d}\rangle(H,H,H)=\sum_{k\mid d}n_k k^3,
\end{equation}
其中 $d\geq1$。由这个等式和上面对 $\langle H,H,H\rangle$ 的定义，经过简单的幂级数运算就能得到式~\eqref{eq:quintic-a}。

遗憾的是，这套论证远不完整，而且似乎预设了 Clemens 猜想。为使一切严格化，我们反过来用式~\eqref{eq:multi-cover} \emph{定义} $n_d$。这样，问题就变成了：$n_d$ 与 $V$ 中的有理曲线究竟有何关系？\hyperref[subsec:instanton-numbers]{第~\ref*{subsec:instanton-numbers} 小节}将证明，当 $d\leq9$ 时，$n_d$ 恰好是五次三维簇上次数为 $d$ 的有理曲线数目。这一点并不像表面看来那样显然。虽然 Clemens 猜想在 $d\leq9$ 时成立，但 $V$ 中有些有理曲线是奇异的。特别地，$V$ 含有 $17{,}601{,}000$ 条次数为 $5$、具有六个结点的有理曲线 \SourceCite{Vainsencher}；要弄清它们如何计入 $n_5$，还须仔细分析。

长期以来，人们以为一旦 Clemens 猜想得到证明，就能立即推出式~\eqref{eq:multi-cover} 中的 $n_d$ 是 $V$ 上次数为 $d$ 的有理曲线数目。然而，R. Pandharipande 最近注意到，即使 Clemens 猜想对 $d=10$ 成立，式~\eqref{eq:multi-cover} 定义的 $n_{10}$ \emph{也不是} $V$ 上十次有理曲线的数目。原因在于，上述六结点五次有理曲线的二重覆盖，所贡献的量超过预期的 $5^3$ 因子。第 \hyperref[chap:gromov-witten]{\ref*{chap:gromov-witten}}、\hyperref[chap:mirror-theorem]{\ref*{chap:mirror-theorem}} 章将对此作严格证明。因此，我们不再把 $n_d$ 称为“次数为 $d$ 的有理曲线的数目”，而采用\hyperref[chap:gromov-witten]{第~\ref*{chap:gromov-witten} 章}中较为中性的名称：\emph{瞬子数}。

对 $n_{10}$ 的讨论体现了\hyperref[chap:gromov-witten]{第~\ref*{chap:gromov-witten} 章}的一个主题：Gromov--Witten 不变量与\IndexLocation{indexloc:0241:39}枚举几何之间存在微妙的相互关系。尽管会遇到这样的困难，\hyperref[chap:gromov-witten]{第~\ref*{chap:gromov-witten} 章}仍将表明，Gromov--Witten 不变量包含大量枚举信息，而镜像对称也推动了枚举代数几何中的一股研究热潮。

本章将依照\hyperref[chap:introduction]{第 1 章}概述的策略，用镜像对称计算五次三维簇 $V$ 的 A 模型关联函数~\eqref{eq:quintic-a}。接下来的几节有一些十分精彩的计算，最终会给出对 $n_d$ 的明确预测。由于我们是用尚不严格的物理理论，来计算枚举意义尚不明朗的数值，这样的讨论难免令人感到不够满意。
\SourcePage{40}
不过，后面的章节将说明，计算中的每一步都可以严格化。特别地，\hyperref[chap:mirror-theorem]{第~\ref*{chap:mirror-theorem} 章}将证明：若用式~\eqref{eq:multi-cover} 以及 $\langle I_{0,3,d}\rangle(H,H,H)$ 的辛几何定义或代数定义来确定 $n_d$，那么镜像对称对 $n_d$ 的预测，便是镜像定理的直接推论。

\section{五次三维簇的镜像}\label{sec:quintic-mirror}

五次三维簇 $V$ 满足 $h^{1,1}(V)=1$，因此它的镜像应满足 $\dim H^1(V^\circ,T_{V^\circ})=1$。所以，五次三维簇的镜像是一个单参数的 Calabi--Yau 流形族。Greene 和 Plesser \SourceCite{GPl} 把它描述为 $\mathbb P^4/G$ 中某个超曲面族的解奇异化。这里，$G$ 是群
\begin{equation}\label{eq:greene-plesser-group}
G=\bigl\{(a_1,\ldots,a_5)\in\mathbb Z_5^5:\textstyle\sum_i a_i\equiv0\pmod5\bigr\}/\mathbb Z_5,
\end{equation}
其中分母中的 $\mathbb Z_5$ 以对角方式嵌入。元素 $g=(a_1,\ldots,a_5)\in G$ 在 $\mathbb P^4$ 上的作用为
\[
g\cdot(x_1,\ldots,x_5)=(\mu^{a_1}x_1,\ldots,\mu^{a_5}x_5),
\]
其中 $\mu=e^{2\pi i/5}$ 是一个本原五次单位根。

以 $\psi$ 为参数，这个超曲面族由方程
\begin{equation}\label{eq:quintic-pencil}
x_1^5+x_2^5+x_3^5+x_4^5+x_5^5+\psi x_1x_2x_3x_4x_5=0
\end{equation}
定义。该方程在 $G$ 的作用下不变，因而在 $\mathbb P^4/G$ 中定义了一个超曲面族。（\SourceCite{CdGP} 用 $-5\psi$ 代替这里的 $\psi$。）商空间 $\mathbb P^4/G$ 是奇异的，式~\eqref{eq:quintic-pencil} 给出的超曲面也继承了这些奇点。不过，只要 $\psi\ne-5\mu^i$，$0\leq i\leq4$，该超曲面就是 $\mathbb P^4/G$ 中的 Calabi--Yau 轨形。更好的是，我们可以同时对所有 $\psi$ 解奇异化，而不破坏 Calabi--Yau 性质；这一点可在局部直接验证。因此，我们把由式~\eqref{eq:quintic-pencil} 得到的解奇异化超曲面族定义为 $V^\circ$。为了强调这是一个单参数簇族，有时记作 $V_\psi^\circ$。

$V_\psi^\circ$ 的这一构造归功于 Greene 和 Plesser，它涉及\hyperref[chap:introduction]{第 1 章}讨论的超共形场论。不过，\hyperref[chap:mirror-constructions]{第 4 章}将说明，Batyrev 对 $V_\psi^\circ$ 给出了纯数学的描述，并将其推广到环面簇中的 Calabi--Yau 超曲面。

如上所述，式~\eqref{eq:quintic-pencil} 定义了奇异簇。由于我们研究的是三维簇，对它的解奇异化并无典范选择。\SourceCite{Morrison2} 给出了一个明确的选择，但还有许多其他选择。幸运的是，不同选择的复模相同，B \IndexLocation{indexloc:1167:40}模型关联函数也相同。我们将用 $V_\psi^\circ$ 指代其中任一选择。（\hyperref[chap:moduli]{第~\ref*{chap:moduli} 章}将看到，这些不同选择反映了镜像的 Kähler 模；在物理学中，它们对应同一个物理理论的不同“相”。）

接下来的任务，是充实第 \ref{sec:three-point} 节讨论的策略，用镜像 $V_\psi^\circ$ 来计算式~\eqref{eq:quintic-a}。第一步是理解有关模空间。对于镜像的复模，注意映射 $(x_1,\ldots,x_5)\mapsto(\mu^{-1}x_1,x_2,\ldots,x_5)$ 诱导同构 $V_\psi^\circ\simeq V_{\mu\psi}^\circ$。\footnote{这一点对于 $\mathbb P^4/G$ 中的超曲面是显然的；还可以证明该同构延伸到解奇异化，见 \SourceCite{Morrison2}。}因此，$\psi^5$ 在 $V_\psi^\circ$ 的复结构模空间上得到良好定义，
\SourcePage{41}
于是
\[
x=\psi^{-5}
\]
是复模的一个局部坐标；选择 $x$ 的原因很快就会明白。此外，可验证 $V_\psi^\circ$ 在 $\psi=-5\mu^i$（$0\leq i\leq4$）以及 $\psi=\infty$ 时发生奇异。用模坐标 $x$ 表示，这意味着复模空间具有边界点 $x=-5^{-5}$ 和 $x=0$。

根据镜像对称，$V^\circ$ 的复模应对应五次三维簇 $V$ 的 Kähler 模。由于 $H$ 生成 $H^2(V,\mathbb Z)\simeq\mathbb Z$，我们可以写 $\omega=tH$。又因为 $H$ 是 Kähler 的，当且仅当 $t$ 位于上半平面时，$\omega=tH$ 是复化 Kähler 类。依照定义~\ref{def:kahler-moduli}，Kähler 空间 $K_{\mathbb C}(V)$ 是复化 Kähler 类对 $H^2(V,\mathbb Z)$ 取商。因此，$q=e^{2\pi it}$ 诱导一个同构
\[
q:K_{\mathbb C}(V)\simeq\Delta^*.
\]
所以，$q$ 是 Kähler 模空间的局部参数，$q=0$ 是一个边界点。一个重要而容易验证的事实是：若 $\ell$ 是 $V$ 中的一条直线，则 $q$ 可写成
\begin{equation}\label{eq:q-kahler}
q=e^{2\pi it}=e^{2\pi i\int_\ell\omega}.
\end{equation}
因此，A 模型关联函数~\eqref{eq:quintic-a} 自然是 Kähler 模局部参数 $q$ 的函数。

假定镜像对称成立，$V$ 的 Kähler 模与 $V^\circ$ 的复模之间便有一个局部同构。参照图式~\eqref{eq:mirror-foliations}，“$V$ 不变、$\omega$ 变化”的截片，局部对应于镜像一侧“$\omega^\circ$ 不变、$V^\circ$ 变化”的截片。然而，关键在于“局部”二字：我们需要给两个模空间各选合适的局部坐标。对 $V$ 的 Kähler 模而言，由于三点函数~\eqref{eq:quintic-a} 是 $q$ 的幂级数，所需坐标显然是 $q$。既然 $q$ 是一个边界点处的局部坐标，镜像一侧也应如此。

在 $V^\circ$ 的复模空间中，可以从两个边界点 $x=0$ 与 $x=-5^{-5}$ 里选择。我们断言 $x=0$ 对应 $q=0$，所以 $x$ 就是需要的局部坐标。理由如下。回到五次三维簇 $V$：与 $H$ 作杯积给出 $\bigoplus_{p=0}^3H^{p,p}(V)$ 上的自同态 $\smile H$，它是最大幂零的，即 $(\smile H)^3\ne0$ 而 $(\smile H)^4=0$。第 \hyperref[chap:hodge-yukawa]{\ref*{chap:hodge-yukawa}}、\hyperref[chap:moduli]{\ref*{chap:moduli}} 章将说明，在镜像对称下，我们预期 $\smile H$ 对应于 $V^\circ$ 的复模空间上、与 $q=0$ 相应的点周围的单值性之对数。因此，镜像对称告诉我们要寻找\emph{最大幺幂单值性}。在上述边界点中，只有 $x=0$ 具有这一性质。\hyperref[chap:hodge-yukawa]{第 5 章}将严格定义最大幺幂单值性。

为了计算 A 模型关联函数~\eqref{eq:quintic-a}，还剩下两项任务：
\begin{itemize}
\item 局部坐标 $q$ 和 $x$ 描述 Kähler 模空间与复模空间中相应的边界点，但镜像对称给出的局部同构\emph{不是} $q=x$。从物理学的观点看，还要考虑量子修正。因此，需要把 $x$ 表示为 $q$ 的函数，反过来也一样。完成后，便得到 $V^\circ$ 的复模空间与 $V$ 的 Kähler 模空间之间的显式局部映射；这通常称为\emph{镜像映射}。

\SourcePage{42}
\item 我们还必须计算 B 模型关联函数并正确地将它归一化，然后用镜像映射把它写成 $q$ 的函数。这样就会得到式~\eqref{eq:quintic-a} 的显式公式。
\end{itemize}
现在先着手第一项任务。

\section{镜像映射}\label{sec:quintic-mirror-map}

理解镜像映射的一种方式是：$q$ 是 $V$ 的 Kähler 模空间上典范选出的坐标。因此，我们希望在 $V^\circ$ 的模空间上找到一个“典范”坐标。\hyperref[chap:moduli]{第~\ref*{chap:moduli} 章}将解释，之所以能这样做，是因为 $x=0$ 处的单值性是\IndexLocation{indexloc:0772:42}最大幺幂的。按照 \SourceCite{Morrison2}，在 $x=0$ 附近存在一个极小整消失圈 $\gamma_0$，它在单值性作用下保持不变。此外，还有一个极小整圈 $\gamma_1$，绕 $x=0$ 的单值性将其变为 $\gamma_1+m\gamma_0$，其中 $m\in\mathbb Z$。对\IndexLocation{indexloc:1181:42}\IndexLocation{indexloc:1195:42}五次三维簇，$m=1$ \SourceCite{Morrison2}。因此，若 $\Omega$ 是 $V^\circ$ 上的全纯三形式，$\int_{\gamma_1}\Omega/\int_{\gamma_0}\Omega$ 在单值性作用下变为它自身加 $1$，从而
\[
\exp\!\left(2\pi i\frac{\int_{\gamma_1}\Omega}{\int_{\gamma_0}\Omega}\right)
\]
是 $x=0$ 邻域中一个得到良好定义的函数。此外，$\gamma_0$ 除符号外唯一；$\gamma_1$ 则在取相同符号以及加上 $\gamma_0$ 的整数倍的意义下唯一。因此，上式在 $V^\circ$ 的模空间上典范地确定。由于 $V^\circ$ 是 Calabi--Yau 簇，它也与 $\Omega$ 的选取无关。于是自然可以断言：镜像映射由 $t=\int_{\gamma_1}\Omega/\int_{\gamma_0}\Omega$ 以及
\begin{equation}\label{eq:mirror-map-periods}
q=\exp\!\left(2\pi i\frac{\int_{\gamma_1}\Omega}{\int_{\gamma_0}\Omega}\right)
\end{equation}
给出。\SourceCite{BCOV2} 给出了支持这一断言的物理学论证；我们还可以从数学上证明，这个镜像映射确实正确，见定理~\ref{thm:lly-quintic-mirror} 或定理~\ref{thm:givental-projective-mirror}。

为了把 $\int_{\gamma_1}\Omega$ 与 $\int_{\gamma_0}\Omega$ 明确表示为 $x$ 的函数，我们沿用 \SourceCite{Morrison1} 的方法。对于 $V^\circ$ 上任意三维圈 $\gamma$，积分 $y=\int_\gamma\Omega$ 称为该族的一个\emph{周期}。众所周知，周期满足 Picard--Fuchs 方程。对 $h^3(V^\circ)=4$ 的三维簇，该方程的形式为
\[
y^{(4)}+f_1y^{(3)}+f_2y''+f_3y'+f_4y=0,
\]
其中 $f_i$ 依赖于模坐标 $x$，导数也对 $x$ 取。一般地，Picard--Fuchs 方程可通过多种方法计算：直接计算周期 \SourceCite{CdGP}；用 Griffiths--Dwork 方法降低极点阶数 \SourceCite{Morrison1,Font,KT1}；利用广义超几何方程 \SourceCite{BvS,HKTY1}；或者，按照 \SourceCite{Givental4} 的一个断言，用环面完全交的显式公式。后三种方法将在\hyperref[chap:hodge-yukawa]{第 5 章}讨论。\hyperref[chap:mirror-theorem]{第~\ref*{chap:mirror-theorem} 章}还会给出整理这些数据的其他方式，以及不同版本镜像定理的证明。第 \ref{sec:quintic-other-methods} 节将先作简要预览。

若把 $\Omega$ 选为三形式
\begin{equation}\label{eq:quintic-residue}
\begin{aligned}
\Omega=\operatorname{Res}\Bigg(\frac{\psi}{x_1^5+\cdots+x_5^5+\psi x_1\cdots x_5}
\Big(&x_1\,dx_2\wedge dx_3\wedge dx_4\wedge dx_5\\
&-x_2\,dx_1\wedge dx_3\wedge dx_4\wedge dx_5\\
&+x_3\,dx_1\wedge dx_2\wedge dx_4\wedge dx_5\\
&-x_4\,dx_1\wedge dx_2\wedge dx_3\wedge dx_5\\
&+x_5\,dx_1\wedge dx_2\wedge dx_3\wedge dx_4\Big)\Bigg),
\end{aligned}
\end{equation}
\SourcePage{43}
那么，\hyperref[chap:hodge-yukawa]{第 5 章}的方法便给出 Picard--Fuchs 方程
\begin{equation}\label{eq:quintic-pf}
\begin{aligned}
0={}&\left(x\frac{d}{dx}\right)^4y
+\frac{2\cdot5^5x}{1+5^5x}\left(x\frac{d}{dx}\right)^3y
+\frac{7\cdot5^4x}{1+5^5x}\left(x\frac{d}{dx}\right)^2y\\
&+\frac{2\cdot5^4x}{1+5^5x}\left(x\frac{d}{dx}\right)y
+\frac{24\cdot5x}{1+5^5x}y.
\end{aligned}
\end{equation}
该方程在 $x=0$ 处有一个正则奇点。还要注意，在另一个边界点 $x=-5^{-5}$ 处也有奇点。

周期 $y_0=\int_{\gamma_0}\Omega$ 与 $y_1=\int_{\gamma_1}\Omega$ 都满足这个方程。$\gamma_0$ 与 $\gamma_1$ 的单值性性质说明，$y_0$ 和 $2\pi i y_1-(\log x)y_0$ 在 $x=0$ 处都是单值的。再用求解式~\eqref{eq:quintic-pf} 的标准方法，可知 $y_0,y_1$ 在下述变换意义下唯一：
\[
y_0\longmapsto b_1y_0,\qquad y_1\longmapsto b_1y_1+b_2y_0,
\]
其中 $b_1,b_2\in\mathbb C$ 且 $b_1\ne0$。这与 $\gamma_0,\gamma_1$ 的唯一性类似，只是这里丢失了整数格。因此，我们可以确定 $q=e^{2\pi i y_1/y_0}$，但仍相差一个常数因子 $c_1=\exp(2\pi i b_2/b_1)$。所得镜像映射为
\begin{equation}\label{eq:q-series-c1}
q=c_1(x-770x^2+\cdots),
\end{equation}
其逆映射为
\begin{equation}\label{eq:x-series-c1}
x=\frac q{c_1}+770\left(\frac q{c_1}\right)^2+\cdots.
\end{equation}
第 \ref{sec:quintic-assemble} 节将利用\IndexLocation{indexloc:1197:43}镜像对称证明 $c_1=-1$。也有其他方法可得到这种归一化；例如，在求解\IndexLocation{indexloc:0618:43}广义\IndexLocation{indexloc:0617:43}超几何函数时采用 Frobenius 方法，这种归一化就是一种自然的选择。\SourceCite{BvS} 提到了这个想法。\hyperref[chap:moduli]{第~\ref*{chap:moduli} 章}一般性地讨论镜像映射时，我们将再回到这一点。

利用广义超几何方程的求解方法，许多公式都能写成闭式。例如，在常数倍意义下，$y_0$ 为
\begin{equation}\label{eq:y0-hypergeometric}
y_0=\sum_{n=0}^{\infty}\frac{(5n)!}{(n!)^5}(-1)^n x^n=1-120x+\cdots;
\end{equation}
若假定 $c_1=-1$，则镜像映射可写为
\begin{equation}\label{eq:q-hypergeometric}
q=-x\exp\!\left(\frac{5}{y_0(x)}\sum_{n=1}^{\infty}\frac{(5n)!}{(n!)^5}
\left[\sum_{j=n+1}^{5n}\frac1j\right](-1)^n x^n\right).
\end{equation}
推导细节见\hyperref[sec:mirror-map]{第~\ref*{sec:mirror-map} 节}；第 \ref{sec:quintic-other-methods} 节还会介绍另一种方法。结合式~\eqref{eq:q-hypergeometric} 与 Dwork 的 $p$ 进超几何级数理论，可以证明 $q$ 是 $x$ 的一个整系数幂级数 \SourceCite{LY2}。

\SourcePage{44}
\section{B 模型关联函数}\label{sec:quintic-b-model}

现在转向计算 $V^\circ$ 上的 Yukawa 耦合。一个关键但初等的观察是：Yukawa 耦合本身也满足一个线性微分方程，而且可以很容易地从 Picard--Fuchs 方程推导出来。依照第 \ref{sec:three-point} 节的式~\eqref{eq:b-yukawa}，有
\begin{equation}\label{eq:b-x-derivatives}
\left\langle x\frac d{dx},x\frac d{dx},x\frac d{dx}\right\rangle
=\int_{V^\circ}\Omega\wedge\Omega''',
\end{equation}
其中撇号表示算子 $x\frac d{dx}$。这里假设 $\Omega$ 由式~\eqref{eq:quintic-residue} 给出。我们也可以把 $\Omega$ 乘以 $x$ 的任意处处非零函数；究竟选哪个倍数，就是\emph{归一化问题}，留待本节稍后处理。正如第 \ref{sec:three-point} 节脚注所说，归一化还会在式~\eqref{eq:b-x-derivatives} 右边引入一个负号；\hyperref[subsec:normalized-yukawa-mirror]{第~\ref*{subsec:normalized-yukawa-mirror} 小节}将对此作解释。

我们断言，Yukawa 耦合 $Y=\int_{V^\circ}\Omega\wedge\Omega'''$ 满足微分方程
\begin{equation}\label{eq:yukawa-ode}
\left(x\frac d{dx}\right)Y=-\frac{5^5x}{1+5^5x}Y.
\end{equation}
为验证这一点，先注意：由 Griffiths 横截性，$\Omega'$ 和 $\Omega''$ 都没有 $(0,3)$ 分量。因此
\begin{equation}\label{eq:omega-zero}
\int_{V^\circ}\Omega\wedge\Omega
=\int_{V^\circ}\Omega\wedge\Omega'
=\int_{V^\circ}\Omega\wedge\Omega''=0.
\end{equation}
将最后一个等式求导两次，得到
\begin{equation}\label{eq:omega-identity}
\int_{V^\circ}\Omega\wedge\Omega^{(4)}
+2\int_{V^\circ}\Omega'\wedge\Omega'''=0.
\end{equation}
另一方面，由 $Y=\int_{V^\circ}\Omega\wedge\Omega'''$ 得
$Y'=\int_{V^\circ}\Omega'\wedge\Omega'''+\int_{V^\circ}\Omega\wedge\Omega^{(4)}$。利用式~\eqref{eq:omega-identity}，可知 $Y'=\frac12\int_{V^\circ}\Omega\wedge\Omega^{(4)}$。而 Picard--Fuchs 方程~\eqref{eq:quintic-pf} 蕴含上同调中的等式
\begin{equation}\label{eq:omega-pf}
0=\Omega^{(4)}+\frac{2\cdot5^5x}{1+5^5x}\Omega'''
+\frac{7\cdot5^4x}{1+5^5x}\Omega''
+\frac{2\cdot5^4x}{1+5^5x}\Omega'
+\frac{24\cdot5x}{1+5^5x}\Omega.
\end{equation}
把 $\Omega$ 与式~\eqref{eq:omega-pf} 右边各项作楔积，再用式~\eqref{eq:omega-zero}，便得到式~\eqref{eq:yukawa-ode}。

式~\eqref{eq:yukawa-ode} 的解为
\begin{equation}\label{eq:yukawa-unnormalized}
Y=\frac{c_2}{1+5^5x},
\end{equation}
其中常数 $c_2$ 将在下一节确定。因此，我们已确定 Yukawa 耦合，但仍相差一个常数因子。可以看到，它在边界点 $x=-5^{-5}$ 处是奇异的，这也符合预期。

但工作尚未完成。式~\eqref{eq:quintic-residue} 给出的 $\Omega$ 可乘以 $x$ 的任意处处非零函数，所以还须归一化 Yukawa 耦合。一种自然的做法，是要求归一化后 $\int_{\gamma_0}\Omega=1$。沿用式~\eqref{eq:quintic-residue} 的 $\Omega$，取周期 $y_0(x)=\int_{\gamma_0}\Omega$，再用 $\Omega/y_0(x)$ 代替 $\Omega$，就得到所需归一化。由式~\eqref{eq:omega-zero} 很容易证明，这种归一化对 Yukawa 耦合
\SourcePage{45}
$Y$ 的作用，是把 $Y$ 换成 $Y/y_0(x)^2$。因此，在这一 $\Omega$ 的选取下，得到\emph{归一化 Yukawa 耦合}
\begin{equation}\label{eq:yukawa-normalized}
\left\langle x\frac d{dx},x\frac d{dx},x\frac d{dx}\right\rangle
=\frac{c_2}{(1+5^5x)y_0(x)^2}.
\end{equation}
注意，上面提到的负号实质上不会影响这个表达式，因为它可以吸收到 $c_2$ 的选取中。据我们所知，这或许也是该负号此前未出现在数学文献中的原因。\hyperref[chap:hodge-yukawa]{第 5 章}将给出计算归一化 Yukawa 耦合的一般方法，并讨论它与 Gauss--Manin 联络之间出乎意料地紧密的联系。

\section{综合起来}\label{sec:quintic-assemble}

现在可以利用前两节的公式，计算五次三维簇的 A 模型关联函数。镜像对称给出同构
\[
H^1(V,\Omega_V^1)\simeq H^1(V^\circ,T_{V^\circ}),
\]
并使 A 模型与 B 模型的关联函数相容。于是，若 $\theta\in H^1(V^\circ,T_{V^\circ})$ 在镜像对称下对应于 $H\in H^1(V,\Omega_V^1)$，依照\hyperref[chap:introduction]{第 1 章}的解释，就应有 $\langle H,H,H\rangle=\langle\theta,\theta,\theta\rangle$。

下面确定 $\theta$。把 $H^1(V,\Omega_V^1)$ 看作 Kähler 模空间的切空间，$H$ 就是 $V$ 的 Kähler 模空间上的一个向量场。若把一般的复化 Kähler 类写为 $tH$，该向量场便对应于 $d/dt$。用 Kähler 模的局部坐标 $q=e^{2\pi it}$ 表示，有 $\frac d{dt}=2\pi iq\frac d{dq}$。在镜像映射 $q=q(x)$ 下，$\frac d{dq}$ 映为 $\frac{dx}{dq}\frac d{dx}$。于是
\[
H=2\pi iq\frac d{dq}\quad\longmapsto\quad
\theta=2\pi iq\frac{dx}{dq}\frac d{dx}.
\]
由 $\langle H,H,H\rangle=\langle\theta,\theta,\theta\rangle$，得到
\begin{equation}\label{eq:quintic-mirror-yukawa}
\begin{aligned}
\langle H,H,H\rangle
&=\left(2\pi iq\frac{dx}{dq}\right)^3
  \left\langle\frac d{dx},\frac d{dx},\frac d{dx}\right\rangle\\
&=\left(2\pi i\frac qx\frac{dx}{dq}\right)^3
  \left\langle x\frac d{dx},x\frac d{dx},x\frac d{dx}\right\rangle.
\end{aligned}
\end{equation}
这一断言就是 A 模型与 B 模型关联函数通过镜像映射相等，对五次三维簇而言，这正是镜像对称。

为使式~\eqref{eq:quintic-mirror-yukawa} 更明确，注意式~\eqref{eq:yukawa-normalized} 已计算出右边的归一化 Yukawa 耦合。把它与镜像映射~\eqref{eq:q-series-c1}、\eqref{eq:x-series-c1} 以及 $y_0(x)$ 的公式~\eqref{eq:y0-hypergeometric} 结合，便有
\[
\begin{aligned}
\langle H,H,H\rangle
&=\frac{c_2}{(1+5^5x)y_0(x)^2}
  \left(2\pi i\frac qx\frac{dx}{dq}\right)^3\\
&=\frac{c_2}{(1+5^5x)(1-120x+\cdots)^2}
 \left(2\pi i(1-770x+\cdots)
 \left(1+1540\frac q{c_1}+\cdots\right)\right)^3.
\end{aligned}
\]
再用式~\eqref{eq:x-series-c1} 全部换成 $q$ 坐标，化简为
\begin{equation}\label{eq:quintic-c-series}
\langle H,H,H\rangle
=\frac{(2\pi i)^3c_2}{
 (1+5^5(q/c_1)+\cdots)(1-240(q/c_1)+\cdots)}
 \left(1+770\frac q{c_1}+\cdots\right)^3.
\end{equation}

\SourcePage{46}
现在确定常数 $c_1$ 和 $c_2$。基本思路是让上述幂级数与我们已知的、式~(2.1) 中给出的公式
\begin{equation}\label{eq:quintic-a-expand}
\langle H,H,H\rangle=5+\sum_{d=1}^{\infty}n_d d^3\frac{q^d}{1-q^d}
=5+2875q+\cdots
\end{equation}
相吻合。首先，由于这个级数的常数项为 $5$，与式~\eqref{eq:quintic-c-series} 比较便得到 $c_2=5/(2\pi i)^3$。其次，我们断言 $c_1=-1$。事实上，采用上述 $c_2$ 后，式~\eqref{eq:quintic-c-series} 的幂级数展开为
\[
\langle H,H,H\rangle=5-2875\frac q{c_1}+\cdots.
\]
经典结果表明，一般五次三维簇包含 $2875$ 条直线，所以式~\eqref{eq:quintic-a-expand} 中 $n_1=2875$。因此，为使两个级数相等，必须有 $c_1=-1$。这并非一厢情愿；\SourceCite{CdGP} 的原始计算确实给出了所需的值 $c_1=-1$。如前所说，\hyperref[chap:moduli]{第~\ref*{chap:moduli} 章}研究镜像映射时，还会给出系统确定 $c_1$ 的方法。

确定 $c_1,c_2$ 后，可以把五次三维簇 $V$ 的镜像对称概括如下。首先，根据式~\eqref{eq:y0-hypergeometric}，
\begin{equation}\label{eq:quintic-y0-final}
y_0(x)=\sum_{n=0}^{\infty}\frac{(5n)!}{(n!)^5}(-1)^nx^n
\end{equation}
是镜像 $V^\circ$ 的 Picard--Fuchs 方程~\eqref{eq:quintic-pf} 唯一的全纯解，唯一性只差一个常数倍。还有一个具有对数奇点的解
\begin{equation}\label{eq:quintic-y1-final}
y_1(x)=y_0(x)\log(-x)
+5\sum_{n=1}^{\infty}\frac{(5n)!}{(n!)^5}
\left[\sum_{j=n+1}^{5n}\frac1j\right](-1)^nx^n.
\end{equation}
依照式~\eqref{eq:q-hypergeometric}，镜像映射为
\begin{equation}\label{eq:quintic-q-final}
q=\exp\!\left(\frac{y_1}{y_0}\right)
=-x\exp\!\left(\frac5{y_0(x)}\sum_{n=1}^{\infty}\frac{(5n)!}{(n!)^5}
\left[\sum_{j=n+1}^{5n}\frac1j\right](-1)^nx^n\right).
\end{equation}
那么，式~\eqref{eq:quintic-mirror-yukawa} 所述五次三维簇的镜像对称断言
\begin{equation}\label{eq:quintic-mirror-identity}
5+\sum_{d=1}^{\infty}n_dd^3\frac{q^d}{1-q^d}
=\frac5{(1+5^5x)y_0(x)^2}\left(\frac qx\frac{dx}{dq}\right)^3.
\end{equation}
请记住：左边是 $V$ 的 A 模型关联函数；右边是 $V^\circ$ 的 B 模型关联函数，即把归一化 Yukawa 耦合视为 $q$ 的函数。

展开式~\eqref{eq:quintic-mirror-identity} 的右边，可得
\begin{equation}\label{eq:quintic-first-coefficients}
5+\sum_{d=1}^{\infty}n_dd^3\frac{q^d}{1-q^d}
=5+2875\frac q{1-q}+609250\cdot2^3\frac{q^2}{1-q^2}+\cdots.
\end{equation}
所以，镜像对称蕴含 $n_2=609250$。继续计算这个级数的其他项，便得到五次三维簇瞬子数 $n_d$ 的下列预测：

\SourcePage{47}
\begin{center}
\begin{tabular}{r|r}
$d$ & $n_d$\\\hline
1 & 2875\\
2 & 609250\\
3 & 317206375\\
4 & 242467530000\\
5 & 229305888887625
\end{tabular}
\end{center}
前四项与已知的枚举结果吻合：$n_1=2875$ 是经典结果（如前所述）；$n_2=609250$ 由 \SourceCite{Katz2} 证明；$n_3=317206375$ 由 \SourceCite{ES1} 证明；$n_4=242467530000$ 由 \SourceCite{Kontsevich2} 证明。\SourceCite{CdGP} 计算了所有 $d\leq10$ 的这些数。

对于一般的 $d$，式~\eqref{eq:quintic-mirror-identity} 断言，式~\eqref{eq:quintic-first-coefficients} 右边级数中出现的 $2875$、$609250$ 等数，正是由五次三维簇的 Gromov--Witten 不变量通过式~\eqref{eq:multi-cover} 定义的瞬子数 $n_d$。\hyperref[subsec:instanton-numbers]{第~\ref*{subsec:instanton-numbers} 小节}将证明：当 $d\leq9$ 时，$n_d$ 的确等于一般五次三维簇上次数为 $d$ 的有理曲线数目。对 $d\geq10$，瞬子数 $n_d$ 仍然存在，但正如已经指出的，它与 $d$ 次有理曲线之间的关系更微妙。

上表中的数暗示 $n_d$ 具有一些非凡的\IndexLocation{indexloc:0630:47}整除性质。容易看出 $n_d$ 是有理数，但尚不知道它们是否总是整数，尽管已计算出的数都是整数。如果假定 $n_d$ 都是整数，那么当 $d$ 不被 $5$ 整除时，可以证明 $5^3\mid n_d$ \SourceCite{LY1}。人们猜想，对\emph{所有} $d$ 都有 $5^3\mid n_d$。

到目前为止，我们还没有谈如何证明式~\eqref{eq:quintic-mirror-identity}。这恰恰是镜像定理的核心内容，也就是下面将要讨论的主题。

\begingroup\linespread{1.22}\selectfont
\section{镜像定理}\label{sec:quintic-other-methods}

前几节的式~\eqref{eq:quintic-mirror-identity} 基于镜像对称的“经典”表述，其首要目标是证明某些关联函数可通过镜像映射相互等同。为了论证这些计算，必须证明一个镜像定理。本节将介绍镜像定理的三个不同版本，它们都蕴含五次三维簇的“经典”镜像对称。我们暂不一般地陈述这些定理，而是考察它们对于五次三维簇 $V$ 及其镜像 $V^\circ$ 的具体形式。完整陈述见第 \hyperref[chap:quantum-cohomology]{\ref*{chap:quantum-cohomology}}、\hyperref[chap:mirror-theorem]{\ref*{chap:mirror-theorem}} 章。

\subsection{Hodge 理论}\label{subsec:quintic-hodge-theory}
第 \ref{sec:quintic-mirror-map} 节构造的镜像映射，把五次三维簇 $V$ 的 Kähler 参数 $q$ 明确表示为镜像 $V^\circ$ 的模参数 $x$ 的函数。因此，我们可以用 $q$ 作为 $V^\circ$ 的模空间的新局部参数。第 \ref{sec:quintic-b-model} 节又确定了归一化 Yukawa 耦合~\eqref{eq:yukawa-normalized}；\hyperref[chap:hodge-yukawa]{第 5 章}将把这一耦合记为 $Y$。前一节的枚举预测，是通过把 $Y$ 写成 $q$ 的函数，再令它等于 A 模型关联函数~\eqref{eq:quintic-a-expand} 得到的。

不过，\hyperref[chap:hodge-yukawa]{第 5 章}将告诉我们，一旦用 $q$ 表示 $Y$，也就完全确定了 $H^3(V^\circ,\mathbb C)$ 上的 Hodge 结构变分。更准确地说，在命题~\ref{prop:normalized-yukawa-pf} 的证明中，我们将看到，在适当的基下，
\SourcePage{48}
$V^\circ$ 的 Gauss--Manin 联络 $\nabla$ 的联络矩阵为
\begin{equation}\label{eq:connection-matrix}
\begin{pmatrix}
0&0&0&0\\
1&0&0&0\\
0&Y&0&0\\
0&0&1&0
\end{pmatrix},
\end{equation}
这里使用 $q$ 作为模参数。

这表明，归一化 B 模型关联函数 $Y$ 并不是物理学所提示的一个奇特函数，而是镜像族的 Hodge 结构不可缺少的一部分。此外，$Y$ 应与 A 模型关联函数对应，这暗示 A 模型一侧也应存在一个 Hodge 结构变分。这是我们第一次看到，镜像对称或可用 Hodge 结构变分来表述。\hyperref[chap:quantum-cohomology]{第~\ref*{chap:quantum-cohomology} 章}定义\emph{A 型 Hodge 结构变分}时，将更充分地发展这一思想。与之相伴的联络称为\emph{A 模型联络}，由量子上同调构造。特别地，对于五次三维簇 $V$，例题~\ref{ex:a-model-yukawa-connection} 将说明：只要把 $Y$ 取为 A 模型关联函数~\eqref{eq:quintic-a-expand}，A 模型联络的联络矩阵也正是式~\eqref{eq:connection-matrix}。

\hyperref[chap:quantum-cohomology]{第~\ref*{chap:quantum-cohomology} 章}将严格定义\emph{数学镜像对}。对五次三维簇及其镜像而言，这意味着 $V$ 的 A 型 Hodge 结构变分，应通过镜像映射同构于镜像族 $V^\circ$ 所产生的 Hodge 结构变分。\hyperref[chap:quantum-cohomology]{第~\ref*{chap:quantum-cohomology} 章}还将证明，这一版本的镜像对称，与式~\eqref{eq:quintic-mirror-identity} 给出的 A 模型、B 模型关联函数相等是等价的。

无论证明镜像对称的经典形式，还是刚讨论的更精细的 Hodge 理论形式，都不是易事。特别是，证明需要引入一些新的研究对象。下面将说明，这些新对象又引出了镜像定理的新版本。

\subsection{Givental 的 $I$ 函数与 $J$ 函数}\label{subsec:givental-ij}

在 Givental 证明镜像定理的方法中，两个上同调值形式函数 $I_{\mathcal V}$ 和 $J_{\mathcal V}$ 起关键作用。先讨论 $J$。这个函数由引力关联数定义；引力关联数可视为 Gromov--Witten 不变量的推广。事实上，我们要先从 Givental 的 $J$ 函数谈起。\hyperref[subsec:cy-threefold-j-function]{第~\ref*{subsec:cy-threefold-j-function} 小节}将看到，对五次三维簇有
\begin{equation}\label{eq:givental-j}
J=e^{(t_0+t_1H)/\hbar}
\left(1+\hbar^{-2}\sum_{d>0}N_dq^d\,d\ell
-2\hbar^{-3}\sum_{d>0}N_dq^d\,\mathrm{pt}\right),
\end{equation}
其中 $q=e^{t_1}$，$\ell$ 是 $V$ 中一条直线的类，$\mathrm{pt}$ 是一个点的类，而且
\begin{equation}\label{eq:nd-gw}
N_d=\langle I_{0,0,d}\rangle=\sum_{k\mid d}n_{d/k}k^{-3}.
\end{equation}
这里的 $n_{d/k}$ 是我们想计算的瞬子数。本章前面提到了 Gromov--Witten 不变量 $\langle I_{0,3,d}\rangle(H,H,H)$。\hyperref[chap:gromov-witten]{第~\ref*{chap:gromov-witten} 章}的 Gromov--Witten 不变量性质蕴含 $\langle I_{0,3,d}\rangle(H,H,H)=d^3\langle I_{0,0,d}\rangle$，所以 $N_d$ 的上述公式与式~\eqref{eq:multi-cover} 等价。

$J$ 函数在\hyperref[subsec:giventals-j-function]{第~\ref*{subsec:giventals-j-function} 小节} Givental 的量子微分算子理论中居于核心地位。对于眼下的目的，式~\eqref{eq:givental-j} 最关键的性质，是 $J$ 包含我们所关注的全部枚举信息。

\SourcePage{49}
函数 $J$ 取值于 $V$ 的上同调环。但若要把它与下面将介绍的 $I_{\mathcal V}$ 比较，我们需要一个取值于 $H^*(\mathbb P^4)$ 的对象。设 $i:V\to\mathbb P^4$ 是包含映射，$i_!$ 是 Gysin 映射，定义 $J_{\mathcal V}=i_!(J)$ 即可。由式~\eqref{eq:givental-j} 容易得到
\begin{equation}\label{eq:givental-jv}
J_{\mathcal V}=e^{(t_0+t_1H)/\hbar}5H
\left(1+\hbar^{-2}\sum_{d>0}N_dq^d\,\frac d5H^2
-2\hbar^{-3}\sum_{d>0}N_dq^d\,\frac15H^3\right).
\end{equation}
（\hyperref[chap:mirror-theorem]{第~\ref*{chap:mirror-theorem} 章}将以略有不同的方式定义 $J_{\mathcal V}$。）这里重要的是：$J_{\mathcal V}$ 唯一确定各个 $n_d$。

Givental 的第二个函数是 $I_{\mathcal V}$，它也取值于上同调。对五次三维簇 $V$，例题~\ref{ex:givental-quintic-mirror-potential} 给出的公式是
\[
I_{\mathcal V}=e^{(t_0+t_1H)/\hbar}5H
\sum_{d=0}^{\infty}e^{dt_1}
\frac{(5H+\hbar)\cdots(5H+5d\hbar)}{
\bigl((H+\hbar)\cdots(H+d\hbar)\bigr)^5}.
\]
虽然 $I_{\mathcal V}$ 取值于 $\mathbb P^4$ 的上同调，但 Givental 的一个观察是：$I_{\mathcal V}$ 由镜像 $V^\circ$ 的周期确定。因此，$I_{\mathcal V}$ 包含镜像的模信息。Givental 版本的镜像定理断言，$I_{\mathcal V}$ 与 $J_{\mathcal V}$ 之间有一种特别好的关系。我们对证明的解释将使用\emph{量子超平面截面原理}；该原理涉及稳定映射模空间上等变上同调的局部化公式。\hyperref[chap:mirror-theorem]{第~\ref*{chap:mirror-theorem} 章}将讨论所有这些内容。

为看清它与前几节的计算有什么关系，令 $t_0=0$、$\hbar=1$，把 $I_{\mathcal V}$ 的上述公式简化为
\begin{equation}\label{eq:iv-simple}
I_{\mathcal V}=e^{t_1H}5H\sum_{d=0}^{\infty}e^{dt_1}
\frac{\prod_{m=1}^{5d}(5H+m)}{\prod_{m=1}^{d}(H+m)^5}.
\end{equation}
可把 $I_{\mathcal V}$ 展开为
\begin{equation}\label{eq:iv-period-expansion}
I_{\mathcal V}=5H\bigl(y_0(t_1)+y_1(t_1)H+y_2(t_1)H^2+y_3(t_1)H^3\bigr).
\end{equation}
作变量替换 $x=-e^{t_1}$ 后，\hyperref[chap:hodge-yukawa]{第 5 章}将证明 $I_{\mathcal V}$ 满足镜像 $V^\circ$ 的 Picard--Fuchs 方程~\eqref{eq:quintic-pf}。实际上，利用 Frobenius 方法，\hyperref[chap:moduli]{第~\ref*{chap:moduli} 章}将证明 $y_0,y_1,y_2,y_3$ 构成解的一组基。而且，$y_0$ 就是式~\eqref{eq:quintic-y0-final} 的级数，$y_1$ 则是式~\eqref{eq:quintic-y1-final}。因此，在 Givental 的表述中，镜像映射~\eqref{eq:quintic-q-final} 自然以 $q=\exp(y_1/y_0)$ 的形式出现。

Givental 版本的镜像定理等价于如下断言：通过镜像映射变换坐标之后，$J_{\mathcal V}=I_{\mathcal V}/y_0$，见定理~\ref{thm:givental-projective-mirror}。由于式~\eqref{eq:nd-gw}、\eqref{eq:givental-jv} 说明 $J_{\mathcal V}$ 包含 $n_d$，而式~\eqref{eq:iv-simple} 明确给出了 $I_{\mathcal V}$，我们可以据此计算任意多项瞬子数 $n_d$。

不过，为使这一组思想首尾相接，还须说明它与关联函数的关系。先注意：第 \ref{sec:quintic-assemble} 节 $V^\circ$ 的归一化 Yukawa 耦合 $Y$，会自然地从 $I_{\mathcal V}$ 产生。为说明原因，考虑 $I_{\mathcal V}/y_0$。只展开到含 $e^{t_1}$ 乘以 $t_1$ 的多项式的项，得到
\[
\begin{aligned}
I_{\mathcal V}/y_0=5H\Big(&1+(t_1+770e^{t_1}+\cdots)H\\
&+\bigl(\tfrac{t_1^2}{2}+770t_1e^{t_1}+575e^{t_1}+\cdots\bigr)H^2\\
&+\bigl(\tfrac{t_1^3}{6}+385t_1^2e^{t_1}+575t_1e^{t_1}-1150e^{t_1}+\cdots\bigr)H^3\Big).
\end{aligned}
\]

\SourcePage{50}
若改用 $q$ 表示，则
\[
I_{\mathcal V}/y_0=5H\Big(1+sH+\bigl(\tfrac{s^2}{2}+575q+\cdots\bigr)H^2
+\bigl(\tfrac{s^3}{6}+575sq-1150q+\cdots\bigr)H^3\Big),
\]
其中 $s=\log q$。于是，若撇号表示 $d/ds=q\,d/dq$，就可算得
\begin{equation}\label{eq:iv-second-derivative}
\begin{aligned}
(I_{\mathcal V}/y_0)''
&=5(1+575q+\cdots)H^3+5(s+575sq+\cdots)H^4\\
&=(5+2875q+\cdots)(H^3+sH^4).
\end{aligned}
\end{equation}
最后一行的第一个因子以 $5+2875q+\cdots$ 开头；至少到 $q$ 的一次项为止，它与归一化 Yukawa 耦合 $Y$ 相符。实际上，这个因子正是 $Y$。这就是 Givental 关于 $I_{\mathcal V}$ 由镜像周期确定的观察。要证明这一点，可用联络矩阵~\eqref{eq:connection-matrix}。它说明，归一化 Picard--Fuchs 方程可写为
\[
\frac{d^2}{ds^2}\left(\frac{d^2y/ds^2}{Y}\right)=0.
\]
容易看出，此方程有解 $f_0,f_1,f_2,f_3$，使得 $f_0=1$、$f_1=s$、$f_2''=Y$ 且 $f_3''=Ys$。式~\eqref{eq:iv-period-expansion} 给出了未归一化的 Picard--Fuchs 方程~\eqref{eq:quintic-pf} 的解 $y_0,y_1,y_2,y_3$。归一化就是除以 $y_0$；所得解为 $y_0/y_0=1$、$y_1/y_0=s$、$y_2/y_0$ 和 $y_3/y_0$。精彩之处在于，$y_2/y_0$ 与 $y_3/y_0$ 基本上就是 $f_2$ 与 $f_3$。严格地说，有恒等式
\begin{equation}\label{eq:period-derivative-identities}
5\frac{d^2}{ds^2}\frac{y_2}{y_0}=Y,
\qquad
5\frac{d^2}{ds^2}\frac{y_3}{y_0}=Ys.
\end{equation}
要证明第一个等式，注意对某些常数 $a,b,c,d$，有 $y_2/y_0=a+bs+cf_2+df_3$，所以 $(y_2/y_0)''=cY+dYs$。由式~\eqref{eq:iv-second-derivative} 可知 $(y_2/y_0)''=1+575q+\cdots$；又因为 $Y=5+2875q+\cdots$，于是 $c=1/5$、$d=0$。第一个等式得证；第二个等式的论证类似。

利用式~\eqref{eq:period-derivative-identities} 可得 $(I_{\mathcal V}/y_0)''=Y(H^3+sH^4)$。命题~\ref{prop:cy-threefold-j-potential} 将类似地证明 $J_{\mathcal V}''=\langle H,H,H\rangle(H^3+sH^4)$。因此，按上述方式理解的定理~\ref{thm:givental-projective-mirror} 蕴含 $\langle H,H,H\rangle=Y$，从而证明式~\eqref{eq:quintic-mirror-identity} 所要求的 A 模型与 B 模型关联函数相等。这就建立了五次三维簇的“经典”镜像对称。

\subsection{Euler 数据与 Gromov--Witten 势}\label{subsec:euler-data}

\SourceCite{LLY} 给出了证明镜像定理的另一种方法。这篇论文包含五次三维簇镜像对称的完整证明，\hyperref[chap:mirror-theorem]{第~\ref*{chap:mirror-theorem} 章}将予讨论。这里先对其中的内容作一个简短预览。

\SourceCite{LLY} 引入了几个值得我们考察的对象。第一个是\emph{Euler 数据}：它是某些射影空间的等变上同调中相互相容的一组元素，能够把稳定映射的模空间“线性化”。有两组特殊的 Euler 数据，记为 $\widehat Q$ 和 $\widehat P$；粗略地说，它们分别编码 A 模型与 B 模型关联函数的信息。

由 Euler 数据可以构造一些上同调值函数，不过这里的值属于等变上同调。例如，五次三维簇的 Euler 数据 $\widehat P$ 给出函数

\SourcePage{51}
\[
HG[\mathcal I(\widehat P)](t_1)
=e^{-pt_1/\hbar}\sum_{d\geq0}
\frac{\prod_{m=0}^{5d}(5p-m\hbar)}{
\prod_{m=1}^{d}\prod_{k=0}^{4}(p-\lambda_k-m\hbar)}e^{dt_1}.
\]
参见定义~\ref{def:euler-data-hypergeometric-function} 后的讨论。这里 $p$ 是等变超平面类，而 $\lambda_k$ 是某些等变上同调类。如果令 $\lambda_k\to0$——\hyperref[chap:localization]{第~\ref*{chap:localization} 章}将称之为“\IndexLocation{indexloc:0945:51}非等变极限”——就得到普通上同调。此时 $p$ 变成 $H$；再令 $\hbar=-1$，便得到式~\eqref{eq:iv-simple} 中的 $I_{\mathcal V}$。特别地，式~\eqref{eq:iv-period-expansion} 变为
\begin{equation}\label{eq:euler-p-limit}
\lim_{\lambda\to0}HG[\mathcal I(\widehat P)](t_1)
=I_{\mathcal V}=5H\bigl(y_0(t_1)+y_1(t_1)H+y_2(t_1)H^2+y_3(t_1)H^3\bigr).
\end{equation}
和前面一样，这使我们能够把镜像映射写为 $q=\exp(y_1/y_0)$。我们用 $\Psi(t_1)$ 表示商 $y_1/y_0$，于是镜像映射为 $q=\exp(\Psi(t_1))$。若记 $q=e^s$，便有 $s=\Psi(t_1)=y_1(t_1)/y_0(t_1)$。

迄今处理的是描述 $V^\circ$ 复模的 Euler 数据 $\widehat P$。还有由 $V$ 的 Gromov--Witten 不变量决定的 Euler 数据 $\widehat Q$。理解它最容易的途径，是借助\emph{Gromov--Witten 势}，其定义为
\begin{equation}\label{eq:gw-potential-quintic}
\Phi(s)=\frac56s^3+\sum_{d=1}^{\infty}N_dq^d
=\frac56s^3+\sum_{d=1}^{\infty}\langle I_{0,0,d}\rangle q^d,
\end{equation}
其中 $N_d$ 由式~\eqref{eq:nd-gw} 给出，且 $q=e^s$。函数 $\Phi(s)$ 与 A 模型关联函数关系很近；事实上，很容易看出式~\eqref{eq:quintic-a-expand} 就是 $d^3\Phi(s)/ds^3$。\hyperref[chap:quantum-cohomology]{第~\ref*{chap:quantum-cohomology} 章}将详细研究 Gromov--Witten 势。

现在考察 Euler 数据 $\widehat Q$。它给出等变上同调值函数 $HG[\mathcal I(\widehat Q)](s)$；\hyperref[sec:quintic-mirror-theorem]{第~\ref*{sec:quintic-mirror-theorem} 节}将证明
\begin{equation}\label{eq:euler-q-limit}
\lim_{\lambda\to0}HG[\mathcal I(\widehat Q)](s)
=5H\left(1+sH+\frac{\Phi'}5H^2+\frac{s\Phi'-2\Phi}5H^3\right),
\end{equation}
其中撇号表示 $d/ds$。（\hyperref[chap:mirror-theorem]{第~\ref*{chap:mirror-theorem} 章}给出的这个公式还包含一个 $\hbar$；与上面一样，这里已令它等于 $-1$。）此式体现了前面所说的：$\widehat Q$ 由 Gromov--Witten 不变量确定。

证明中的关键一步，是说明作镜像映射给出的代换 $s=\Psi(t_1)$ 后，式~\eqref{eq:euler-q-limit} 的极限与式~\eqref{eq:euler-p-limit} 的极限相同，只差一个 $y_0$ 因子。更准确地说，推论~\ref{cor:quintic-hypergeometric-mirror-equality} 将证明
\[
\lim_{\lambda\to0}HG[\mathcal I(\widehat Q)](\Psi(t_1))
=\frac1{y_0}\lim_{\lambda\to0}HG[\mathcal I(\widehat P)](t_1).
\]
把它与式~\eqref{eq:euler-q-limit}、\eqref{eq:euler-p-limit} 结合并比较系数，很容易得到恒等式
\begin{equation}\label{eq:lly-potential}
\Phi(\Psi(t_1))=\frac52\left(\frac{y_1y_2}{y_0y_0}-\frac{y_3}{y_0}\right).
\end{equation}
这就是 \SourceCite{LLY} 所给出的五次三维簇\IndexLocation{indexloc:0864:51}\IndexLocation{indexloc:0867:51}镜像定理。

由于 $\Phi(s)$ 是式~\eqref{eq:gw-potential-quintic} 的 Gromov--Witten 势，公式~\eqref{eq:lly-potential} 使我们能计算任意 $d$ 的 $n_d$。但它与三点函数有什么关系？我们已指出，$d^3\Phi/ds^3$ 就是 A 模型关联函数。五次三维簇的镜像对称断言，经由镜像映射，这个函数等于
\SourcePage{52}
B 模型关联函数 $Y$。因此，只要证明
\begin{equation}\label{eq:lly-yukawa}
\frac{d^3}{ds^3}\frac52\left(\frac{y_1y_2}{y_0y_0}-\frac{y_3}{y_0}\right)=Y,
\end{equation}
式~\eqref{eq:lly-potential} 就会推出式~\eqref{eq:quintic-mirror-identity} 所述的镜像对称表述。这里，与前面一样，$y_i$ 经由镜像映射被看作 $s$ 的函数。

为证明式~\eqref{eq:lly-yukawa}，先用式~\eqref{eq:period-derivative-identities} 得到
\[
\frac{d^2}{ds^2}\frac52\left(\frac{y_1y_2}{y_0y_0}-\frac{y_3}{y_0}\right)
=\frac{d^2}{ds^2}\frac52\left(s\frac{y_2}{y_0}-\frac{y_3}{y_0}\right)
=5\frac d{ds}\frac{y_2}{y_0}.
\]
这个直接的证明留给读者。然后再求一次导数，并再次使用式~\eqref{eq:period-derivative-identities}，式~\eqref{eq:lly-yukawa} 立即得证。至此，五次三维簇的镜像对称证明便完成了；当然，从这里到\hyperref[chap:mirror-theorem]{第~\ref*{chap:mirror-theorem} 章}之间，仍有大量实质性的工作要做。

请注意，在这条证明镜像定理的路线中，归一化 Yukawa 耦合 $Y$ 直到计算的最后才出现；此前所有步骤都涉及新的研究对象。Givental 的方法也是如此，只是中间对象略有不同。因此，在第 \hyperref[chap:toric]{\ref*{chap:toric}}--\hyperref[chap:quantum-cohomology]{\ref*{chap:quantum-cohomology}} 章建立表述镜像对称所需的数学时，读者应记住：要在\hyperref[chap:mirror-theorem]{第~\ref*{chap:mirror-theorem} 章}证明镜像对称，还需要一些根本性的新思想。这会花费大量功夫，但最终结果之奇妙，足以回报这些努力。

\par\endgroup
